\answer{ex:04:1}{(а)~$\approx1.1\cdot10^{11}$ бит/Дж. (б)~$\log_2\sqrt{10}\approx1.7$ бита против $81$: в $\approx50$ раз.}
\answer{ex:04:2}{$r^*=(P_0/E_{\text{имп}})\ln\bigl(1/(r^*\Delta t)\bigr)$, $P_0/E_{\text{имп}}=1.16\,\text{с}^{-1}$: $r^*\approx6$ импульсов в секунду, $7.4\cdot10^{10}$ бит/Дж.}
\answer{ex:04:3}{$\mathcal I=\frac{N}{1+(N-1)c}=9.9$ (предел $10$) против $\mathcal I=\frac{N}{1-c}=1111$: в $\approx110$ раз; корреляция вредна, только когда похожа на сигнал.}
\answer{ex:04:4}{$\theta\approx68.7$ и $19.9$ (в единицах $w$), $m=19$ и $m=10$: быстрому получателю нужно почти вдвое меньше совпадающих входов, хотя его средний вход в пять раз меньше.}
\answer{ex:04:5}{$U\tau_{\text{вос}}\ge0.725\,\text{с}$ и $\tau_{\text{вос}}(1-2U)\le0.05\,\text{с}$, что требует $U\ge0.483$; наименьшее $\tau_{\text{вос}}=0.725\,\text{с}$ при $U=1$ ($1.45\,\text{с}$ при $U=0.5$).}
\answer{ex:04:6}{(а)~$\theta\,e/(e-1)\approx1.58\,\theta$ при любых $\tau$ и $\theta$. (б)~$14$ импульсов.}
\answer{ex:04:7}{(а)~$1.62$ бита на слово: $0.77$ --- число импульсов, $0.84$ --- их расположение. (б)~$1.13$ и $1.59$ бита: по $30$ словам оценка теряет почти треть.}
